\documentclass[12pt]{article}
\usepackage[margin=1in]{geometry}
\usepackage[utf8]{inputenc}
\usepackage{latexsym, amsfonts, enumitem, graphicx, environ,enumitem, fancyhdr, color, amssymb, amsthm, amsmath, mathtools, tikz, nicematrix}
\usetikzlibrary{patterns}
\usetikzlibrary{decorations.pathreplacing}
\pagestyle{fancy}
\setlength{\headheight}{65pt}
\newenvironment{problem}[2][Problem]{\begin{trivlist}
    \item[\hskip \labelsep {\bfseries #1}\hskip \labelsep {\bfseries #2.}]}{\end{trivlist}}
\DeclareMathOperator{\Ima}{Im}

\usepackage{hyperref}
\usepackage[capitalize]{cleveref}

\newtheorem{thm}{Theorem}[section]
\newtheorem{lem}[thm]{Lemma}
\newtheorem{cl}[thm]{Claim}
\newtheorem{prop}[thm]{Proposition}
\newtheorem{cor}[thm]{Corollary}
\newtheorem{conj}[thm]{Conjecture}
\newtheorem{obstruction}[thm]{Obstruction}

\usepackage{mdframed}

\theoremstyle{definition}
\newtheorem{definition}[thm]{Definition}
\newtheorem{remark}[thm]{Remark}
\newtheorem{question}{Question}
\newtheorem{example}[thm]{Example}
\newtheorem{exercise}[thm]{Exercise}
\newtheorem{properties}[thm]{Properties}

\usepackage{tcolorbox}
\tcbuselibrary{theorems,skins,breakable}

\tcolorboxenvironment{example}{
enhanced jigsaw,colframe=cyan,interior hidden,
breakable,%before skip=10pt,after skip=10pt
}

\tcolorboxenvironment{thm}{
enhanced jigsaw,colframe=red,interior hidden,
breakable,%before skip=10pt,after skip=10pt
}

\tcolorboxenvironment{prop}{
enhanced jigsaw,colframe=red,interior hidden,
breakable,%before skip=10pt,after skip=10pt
}

\tcolorboxenvironment{properties}{
enhanced jigsaw,colframe=red,interior hidden,
breakable,%before skip=10pt,after skip=10pt
}

\tcolorboxenvironment{lem}{
enhanced jigsaw,colframe=red,interior hidden,
breakable,%before skip=10pt,after skip=10pt
}
\tcolorboxenvironment{cor}{
enhanced jigsaw,colframe=red,interior hidden,
breakable,%before skip=10pt,after skip=10pt
}
\tcolorboxenvironment{obstruction}{
enhanced jigsaw,colframe=red,interior hidden,
breakable,%before skip=10pt,after skip=10pt
}

\tcolorboxenvironment{proof}{
enhanced jigsaw, frame hidden,%interior hidden,
breakable,%before skip=10pt,after skip=10pt
}

\lhead{Avi Chetlin \\ Assignment 01 \\ 02 September 2026}
\rhead{Dr. Michael Martens \\ PHYS 313 Thermodynamics \\ Fall 2026}

\begin{document}
\begin{problem}{1}
  Suppose you have a 1.0L container of pure water, where each molecule is distinguishable.
  Pour the liter of distinguishable water into the ocean and allow complete mixing to occur.
  Then, scoop up 1.0L of water from the ocean.

  Recall that
  \begin{itemize}
    \item Volume of earths oceans = \( 1.35 \times 10^9 \text{km}^3 \),
    \item Density of water = \( 1.0 \)g cm\(^{-3} \),
    \item Molar mass of water = \( 18.0 \)g mol\(^{-1}\),
          \item Avogadro's number = \( 6.022 \times 10^{23} \) molecules/mol.
  \end{itemize}

  \begin{enumerate}[label=(\alph*)]
    \item How many molecules of water are there in a 1.0L container of pure water?
    \item How many molecules of water are there in the Earth's oceans?
    \item Estimate how many of the original, distinguishable water molecules you expect to find in your sample.
  \end{enumerate}
\end{problem}

\begin{problem}{2}
  Cars on a highway travel at different speeds, with the probability of a car traveling between speeds \( v \) and \( v + dv \) is
  \begin{equation}
  \label{eq:1}
  \rho(v) dv = \alpha v \exp(-v/v_0) dv,
  \end{equation}
  where \( \rho(v) \) is the probability density of the speed.

  \begin{enumerate}[label=(\alph*)]
    \item Find the value of the constant \(\alpha\) in terms of \( v_0 \) such that the total probability is
          \begin{equation}
          \label{eq:2}
          \int\limits_0^{\infty}\rho(v) dv = 1.
          \end{equation}
    \item What is the mean value of the speed of the cars?
    \item What is the most likely speed of a car?
    \item How many cars are traveling more than twice the mean speed?
  \end{enumerate}
\end{problem}

\begin{problem}{3}
  Consider a system with 6 identical marbles placed randomly into a box that has 12 spaces available, each holding at most one marble, and divides equally in half.
  Each possible arrangement will be considered a microstate of the system.
  The macrostate of the system will be defined as the number of marbles, \( n_L \), in the left half.
  The marbles are initially in a state with \( n_L = 6 \).

  \begin{enumerate}[label=(\alph*)]
    \item What is the multiplicity of this system if the marbles can move around within their own half, but cannot cross between the left and ride sides?
    \item The constraint is now removed, and the marbles may move freely between the left and right halves.
          The multiplicity of states for the left is \( g_L(n_L) \), similarly for the right side \( g_R(n_R) \).
          Complete the table, recalling that \( n_L+n_R = 6 \).
          \begin{center}

          \begin{tabular}{|c|c|c|c|}
            \hline
            \( n_{L} \) & \( g_L(n_L) \) & \( g_R(n_R) \) & \( g_Lg_R \) \\
            \hline
            \( 0 \) & & & \\
            \hline
            \( 1 \) & & & \\
            \hline
            \( 2 \) & & & \\
            \hline
            \( 3 \) & & & \\
            \hline
            \( 4 \) & & & \\
            \hline
            \( 5 \) & & & \\
            \hline
            \( 6 \) & & & \\
            \hline
          \end{tabular}
          \end{center}
    \item Find the total number of microstates by adding the \( g_Lg_R \) column.
          Does this agree with the number of arrangements of \( 6 \) marbles in \( 12 \) spaces?
    \item What is the most probable macrostate?
          What is the probability the system is in a macrostate with \( n_L = 3 \)?
    \item In the initial state, \( n_L = 6 \) and \( n_R = 0 \).
          What are the entropies \( \sigma_L \) and \( \sigma_R \)?
    \item After the marbles are allowed to move between the halves, what is the total entropy \(\sigma_S\)?
          Note that \( \sigma_S > \sigma_L + \sigma_R \).
  \end{enumerate}
\end{problem}
\end{document}
